\clearpage
\section{Relaciones y cardinalidades}

La tabla resume todas las relaciones del diagrama final. La lectura
\texttt{1:N} se realiza desde la primera entidad hacia la segunda.

\begin{center}
\scriptsize
\renewcommand{\arraystretch}{0.94}
\begin{tabularx}{\textwidth}{@{}>{\bfseries\color{unamblue}}p{0.20\textwidth}p{0.25\textwidth}p{0.10\textwidth}X@{}}
\toprule
Relación & Entidades & Card. & Restricción semántica \\
\midrule
Poseer & Cliente--Tarjeta & 1:1 & Cada tarjeta pertenece a un cliente y cada
cliente posee una sola tarjeta activa. \\
Registrar & Tarjeta--Partida & 1:N & Una tarjeta registra muchas partidas;
cada partida queda en una sola tarjeta. \\
Ejecutar & Máquina--Partida & 1:N & Una máquina ejecuta muchas partidas;
cada partida ocurre en una máquina. \\
Corresponder & Juego--Máquina & 1:N & Un juego puede estar instalado en
varias máquinas; cada máquina corresponde a un juego. \\
Pertenecer & Sucursal--Máquina & 1:N & Una sucursal alberga varias máquinas;
cada máquina pertenece a una sucursal. \\
Ingresar & Tarjeta--Visita & 1:N & Una tarjeta origina muchas visitas; cada
visita se registra con una tarjeta. \\
Proceder en & Sucursal--Visita & 1:N & Una sucursal recibe muchas visitas;
cada visita ocurre en una sucursal. \\
Recibir & Tarjeta--Recarga & 1:N & Una tarjeta recibe varias recargas; cada
recarga se abona a una tarjeta. \\
Efectuar & Sucursal--Recarga & 1:N & Una sucursal efectúa varias recargas;
cada recarga ocurre en una sucursal. \\
Reflejar & Cajero--Recarga & 1:N & Un cajero refleja varias recargas; cada
recarga la atiende un cajero. \\
Vender & Cajero--Tarjeta & 1:N & Un cajero vende varias tarjetas; cada tarjeta
es vendida por un cajero. \\
Clasificar & Tipo de juego--Juego & 1:N & Un tipo clasifica varios juegos;
cada juego se asocia, como máximo, con un tipo. \\
Disponer & Sucursal--Premio & N:M & Una sucursal dispone de varios premios y
un premio puede ofrecerse en varias sucursales. La relación conserva
\texttt{cantidadDisponible}. \\
Hacer & Cliente--Canje & 1:N & Un cliente hace varios canjes; cada canje
corresponde a un cliente. \\
Otorgar & Premio--Canje & 1:N & Un tipo de premio puede aparecer en varios
canjes; cada canje otorga un premio. \\
Saldar & Sucursal--Canje & 1:N & Una sucursal salda varios canjes; cada canje
se atiende en una sucursal. \\
Procesar & Encargado de premios--Canje & 1:N & Un encargado procesa varios
canjes; cada canje es procesado por un encargado. \\
Trabajar & Sucursal--Empleado & 1:N & Una sucursal concentra varios empleados;
cada empleado trabaja en una sucursal. \\
Administrar & Sucursal--Gerente & 1:1 & Cada sucursal tiene un gerente y cada
gerente administra una sucursal. \\
Realizar & Técnico--Mantenimiento & 1:N & Un técnico realiza varios
mantenimientos; cada mantenimiento tiene un técnico responsable. \\
Tener & Máquina--Mantenimiento & 1:N & Una máquina tiene varios registros de
mantenimiento; cada mantenimiento corresponde a una máquina. \\
\bottomrule
\end{tabularx}
\end{center}

Las entidades transaccionales \texttt{Partida}, \texttt{Visita},
\texttt{Recarga}, \texttt{Canje} y \texttt{Mantenimiento} permiten conservar
el historial sin sobrescribir operaciones previas. La relación N:M
\texttt{Disponer} separa la descripción del premio de su existencia local.
